\begin{table}[htbp]
\centering
\caption{Event Study: Credit to Private Sector (\% GDP) around Kenya's Interest Rate Cap}
\label{tab:eventstudy}
\begin{tabular}{lccc}
\hline\hline
  Year & $\hat{\beta}$ (Kenya diff.) & SE & 95\% CI \\
\hline
  2012 & -6.480 & 2.318 & [-11.024, -1.936] \\
  2013 & -5.567 & 1.229 & [-7.975, -3.158] \\
  2014 & -0.329 & 0.945 & [-2.181, 1.523] \\
  2015 (base) & 0 & --- & --- \\[2pt]
\hline
  2016 (cap intro) & -0.596 & 0.288 & [-1.161, -0.032] \\
  2017 & -2.806 & 0.488 & [-3.762, -1.849] \\
  2018 & -4.953 & 0.857 & [-6.633, -3.273] \\
  2019 (repeal yr) & -5.375 & 0.920 & [-7.178, -3.571] \\
\hdashline
  2020 & -5.324 & 2.211 & [-9.659, -0.990] \\
  2021+ & -6.596 & 0.815 & [-8.193, -4.999] \\
\hline\hline
\multicolumn{4}{p{0.7\textwidth}}{\footnotesize \textit{Notes:} Kenya--control country differential in credit/GDP (\%), relative to 2015 baseline. Control countries: Uganda, Tanzania, Rwanda. Country and year fixed effects. Reference year is 2015 ($\hat{\beta}=-1$, normalized to zero). The cap was introduced in September 2016; full cap years are 2017--2019. The repeal was signed November 2019; post-repeal years begin 2020. Heteroscedasticity-robust standard errors.}
\end{tabular}
\end{table}
